\documentclass[12pt]{report}

\usepackage[a4paper,margin=1in]{geometry}
\usepackage[T1]{fontenc}
\usepackage[utf8]{inputenc}
\usepackage{lmodern}
\usepackage{graphicx}
\usepackage{booktabs}
\usepackage{tabularx}
\usepackage{array}
\usepackage{ragged2e}
\usepackage{enumitem}
\usepackage{listings}
\usepackage{xcolor}
\usepackage{float}
\usepackage[hidelinks]{hyperref}
\usepackage{caption}

\captionsetup{font=small}
\setlist[itemize]{leftmargin=*,itemsep=0.35em}
\setlist[enumerate]{leftmargin=*,itemsep=0.35em}
\setlength{\parskip}{0.6em}
\setlength{\parindent}{0pt}
\renewcommand{\arraystretch}{1.25}

\definecolor{codebg}{RGB}{248,248,248}
\definecolor{placeholdergray}{RGB}{110,110,110}

\lstdefinestyle{reportcode}{
  basicstyle=\ttfamily\small,
  backgroundcolor=\color{codebg},
  frame=single,
  breaklines=true,
  columns=fullflexible,
  keepspaces=true
}

\newcolumntype{Y}{>{\RaggedRight\arraybackslash}X}
\newcommand{\placeholder}[1]{\textcolor{placeholdergray}{\textit{#1}}}
\newcommand{\evidencebox}[2]{
  \begin{figure}[H]
    \centering
    \includegraphics[width=0.88\textwidth]{#1}
    \caption{#2}
  \end{figure}
}

\newcommand{\evidenceplaceholder}[2]{
  \begin{figure}[H]
    \centering
    \fbox{
      \parbox{0.88\textwidth}{
        \centering
        \vspace{1.2in}
        \placeholder{#1}
        \vspace{1.2in}
      }
    }
    \caption{#2}
  \end{figure}
}

\newcommand{\lessonchapter}[3]{
  \chapter{Lesson #1: #2}

  \section{Part 1: Goal and Vulnerability Summary}
  \placeholder{State the vulnerability type, main affected component, security impact, and the expected safe behavior for #2.}

  \section{Part 2: Why This Works / Root Cause}
  \placeholder{Explain the security mistake that makes this #3 issue possible. Keep this conceptual and focused on root cause rather than symptoms.}

  \section{Part 3: Environment and Setup}
  \begin{itemize}
    \item DVSA Website URL / API endpoint: https://shfz8h7a28.execute-api.us-east-1.amazonaws.com/Stage
    \item AWS services involved: \placeholder{For example API Gateway, Lambda, Cognito, S3, DynamoDB, SQS, IAM.}
    \item Relevant function, role, or log group: \placeholder{Name the exact resources used in this lesson.}
    \item User or test context: \placeholder{Attacker, victim, admin, or anonymous context as applicable.}
    \item Tools used: \placeholder{Browser DevTools, curl, jq, AWS Console, AWS CLI, Policy Simulator, CloudWatch, etc.}
  \end{itemize}

  \section{Part 4: Reproduction Steps}
  \begin{enumerate}
    \item \placeholder{Describe the baseline or normal behavior first.}
    \item \placeholder{List the exact steps that trigger the vulnerability.}
    \item \placeholder{Include the request, timing, role, or payload details needed for reproducibility.}
    \item \placeholder{Record the observable result that proves the issue.}
  \end{enumerate}

  \section{Part 5: Evidence and Proof}
  \begin{itemize}
    \item \placeholder{Attach the strongest pre-fix evidence for this lesson.}
    \item \placeholder{Use only evidence that directly proves the claim.}
    \item \placeholder{Redact real tokens, secrets, and temporary credentials before submission.}
  \end{itemize}
  \evidenceplaceholder{Replace with pre-fix evidence for Lesson #1 (#2).}{Pre-fix evidence for Lesson #1: #2}

  \section{Part 6: Fix Strategy / Probable Mitigation}
  \begin{itemize}
    \item Fix location: \placeholder{Identify the backend, IAM, cloud configuration, or workflow layer that must change.}
    \item Security property restored: \placeholder{For example authentication integrity, authorization, least privilege, or safe input handling.}
    \item Why this addresses root cause: \placeholder{Explain why the fix is real and not cosmetic.}
  \end{itemize}

  \section{Part 7: Code / Config Changes}
  \placeholder{Paste the real code diff, policy change, CloudFormation change, or configuration update below.}

  \begin{figure}[H]
    \centering
    \fbox{
      \parbox{0.88\textwidth}{
        \vspace{0.4em}
        \placeholder{Replace this box with the actual diff, code snippet, policy JSON, or configuration change.}
        \vspace{2.0in}
      }
    }
    \caption{Replace with the concrete code or configuration change for Lesson #1.}
  \end{figure}

  \section{Part 8: Verification After Fix}
  \begin{enumerate}
    \item \placeholder{Repeat the same test that proved the vulnerability.}
    \item \placeholder{Show that the exploit no longer succeeds.}
    \item \placeholder{Also show that normal intended behavior still works.}
  \end{enumerate}
  \evidenceplaceholder{Replace with post-fix verification evidence for Lesson #1 (#2).}{Post-fix verification evidence for Lesson #1: #2}

  \section[Part 9: Structured Operation and Security Analysis]{Part 9: Structured Operation and\\Security Analysis}
  \subsection*{Table A}
  \begin{table}[H]
    \small
    \centering
    \begin{tabularx}{\textwidth}{|Y|Y|Y|Y|}
      \hline
      \textbf{Intended rule(s)} & \textbf{Artifacts used to infer the rule} & \textbf{Normal evidence} & \textbf{Exploit evidence} \\
      \hline
      \placeholder{State the intended security or workflow rule.} &
      \placeholder{Identify the source used to infer the rule: UI flow, code path, IAM policy, logs, helper guide, etc.} &
      \placeholder{Show normal behavior when the rule is respected.} &
      \placeholder{Show behavior proving the rule was violated.} \\
      \hline
    \end{tabularx}
  \end{table}

  \subsection*{Table B}
  \begin{table}[H]
    \footnotesize
    \centering
    \begin{tabularx}{\textwidth}{|Y|Y|Y|Y|Y|}
      \hline
      \textbf{Why the exploit is a deviation} & \textbf{Deviation class} & \textbf{Fix applied} & \textbf{Post-fix verification} & \textbf{Optional latency / timing note} \\
      \hline
      \placeholder{Explain why the observed result violates the intended rule.} &
      \placeholder{Examples: authentication failure, access control failure, unsafe input handling, logic flaw, DoS, excessive privilege.} &
      \placeholder{Summarize the real corrective action.} &
      \placeholder{Describe the evidence that the deviation is closed.} &
      \placeholder{Use this only when timing or latency matters.} \\
      \hline
    \end{tabularx}
  \end{table}

  \section{Part 10: Takeaway / Lessons Learned}
  \placeholder{End with a short security lesson that ties the #3 weakness back to serverless design, AWS service interaction, and secure engineering practice.}

  \clearpage
}

% Fill in these fields before compiling the final report.
\newcommand{\CourseCode}{ICS-344}
\newcommand{\CourseTitle}{Information Security}
\newcommand{\TeamNo}{2}
\newcommand{\CourseSection}{Sections: 1 \& 2}
\newcommand{\CourseTerm}{252}
\newcommand{\ProjectTitle}{DVSA Vulnerability Discovery and Remediation}
\newcommand{\StudentNames}{Husain Al Muallim$^1$, Ali Al Qatari$^2$, Abdulwahed Zuhair Adi$^2$}
\newcommand{\StudentIDs}{202163730, 202279780, 202256380}
\newcommand{\DVSAWebsiteURL}{http://dvsa-website-145292399251-us-east-1.s3-website.us-east-1.amazonaws.com}
\newcommand{\AWSRegion}{us-east-1}
\newcommand{\SubmissionDate}{End of Sunday, 3$^{rd}$ of May, 2026}

\begin{document}

\hypersetup{pageanchor=false}
\begin{titlepage}
  \centering
  {\Large King Fahd University of Petroleum and Minerals\par}
  \vspace{1.2cm}
  {\Large \CourseCode\ - \CourseTitle\par}
  {\large Team Number: \TeamNo\par}
  {\large \CourseSection\par}
  \vspace{1.4cm}
  {\huge \bfseries \ProjectTitle\par}
  \vspace{1.4cm}

  \begin{tabularx}{\textwidth}{@{}rY@{}}
    Term: & \CourseTerm \\
    Student Name(s): & \StudentNames \\
    Student ID(s): & \StudentIDs \\
    AWS Region: & \AWSRegion \\
    Date: & \SubmissionDate \\
  \end{tabularx}

  \vspace{0.6cm}
  {\small\textbf{DVSA Website URL:} \url{\DVSAWebsiteURL}\par}

  \vfill
  {\large Project Report\par}
\end{titlepage}

\clearpage
\hypersetup{pageanchor=true}
\pagenumbering{roman}
\tableofcontents
\clearpage

\chapter*{Project Overview}
\addcontentsline{toc}{chapter}{Project Overview}

\section*{Scope}
This project involves deploying the Damn Vulnerable Serverless Application (DVSA) into a non-production AWS account, systematically exploiting 10 official vulnerability lessons, applying remediations, and verifying fixes. All work was performed in a controlled environment with no impact on production systems.

\section*{Deployment Summary}
\begin{itemize}
  \item DVSA Website URL: \DVSAWebsiteURL
  \item AWS region: \AWSRegion
  \item CloudFormation stack: \texttt{serverlessrepo-OWASP-DVSA}
  \item Main services used: S3, API Gateway, Lambda, DynamoDB, SQS, Cognito, IAM, CloudWatch, SES, CloudTrail
\end{itemize}

\section*{Architecture and Request Pipeline}
The DVSA follows a standard serverless architecture. The browser loads the static frontend from an S3 bucket. User actions trigger HTTP requests to API Gateway, which routes them to backend Lambda functions (e.g., \texttt{DVSA-ORDER-MANAGER}). These functions interact with DynamoDB for persistence, SQS for async processing, Cognito for authentication, and S3 for file storage. Responses flow back through API Gateway to the client.

\section*{Methodology and Evidence Sources}
\begin{itemize}
  \item Environment validated by confirming DVSA website loads, user registration works, and orders can be placed end-to-end.
  \item Tools used: \texttt{curl}, Python \texttt{requests}, Browser DevTools, AWS Console, AWS CLI, IAM Policy Simulator, and CloudWatch Logs.
  \item Evidence captured via screenshots and terminal output; real tokens and AWS keys redacted before commit.
  \item Fixes verified by re-running the same exploit and confirming failure, then confirming normal workflow still succeeds.
\end{itemize}

\clearpage
\pagenumbering{arabic}

\chapter{Lesson 01: Event Injection}

\section{Part 1: Goal and Vulnerability Summary}
The objective of this lesson is to successfully exploit an Event Injection flaw to achieve Remote Code Execution (RCE) on the backend cloud server, and subsequently remediate the vulnerability. The vulnerability lies within the \texttt{DVSA-ORDER-MANAGER} AWS Lambda function, which acts as the backend for the application's order system. When a user submits an HTTP request, the API Gateway passes the user's JSON payload directly to this Lambda function. Instead of safely parsing this text, the developer used an outdated and inherently insecure Node.js library called \texttt{node-serialize} to deserialize both the request body and the HTTP headers. Because the backend blindly trusts the incoming user input and processes it with this vulnerable library, an attacker can craft a malicious JSON payload containing executable code, bypassing normal application logic.

\section{Part 2: Why This Works / Root Cause}
The root cause is insecure deserialization. The \texttt{DVSA-ORDER-MANAGER} Lambda function processes incoming HTTP requests using \texttt{serialize.unserialize(event.body)}. Under the hood, the \texttt{node-serialize} package uses JavaScript's \texttt{eval()} function to execute code if it detects a specifically formatted string (starting with \texttt{\_\$\$ND\_FUNC\$\$\_}). Because the backend does not validate or sanitize the incoming JSON payload before passing it to this function, the attacker's injected JavaScript payload is executed with the privileges of the Lambda function.

\section{Part 3: Environment and Setup}
\begin{itemize}
  \item API endpoint:\par
  \url{https://shfz8h7a28.execute-api.us-east-1.amazonaws.com/Stage}
  \item AWS services involved: API Gateway, AWS Lambda, CloudWatch Logs.
  \item Relevant function, role, or log group: Lambda function \nolinkurl{DVSA-ORDER-MANAGER} and its CloudWatch Log Group.
  \item User or test context: Unauthenticated external attacker.
  \item Tools used: Mac Terminal, \texttt{curl} for sending HTTP POST requests, and AWS Management Console (CloudWatch).
\end{itemize}

\section{Part 4: Reproduction Steps}
\begin{enumerate}
  \item The baseline behavior of the \texttt{/order} endpoint is to accept a valid JSON payload containing order details and process it.
  \item To trigger the vulnerability, a crafted JSON payload is sent to the \texttt{/order} API endpoint using \texttt{curl}.
  \item The payload contains the malicious \texttt{\_\$\$ND\_FUNC\$\$\_} tag, followed by a self-executing JavaScript function that writes a file to the Lambda's \texttt{/tmp} directory, reads it, and prints the contents to the console error log.
  \item The following command was executed:
\begin{lstlisting}[style=reportcode]
curl -X POST "https://shfz8h7a28.execute-api.us-east-1.amazonaws.com/Stage/order" \
 -H "Content-Type: application/json" \
 -d '{"action": "_$$ND_FUNC$$_function(){ var fs = require(\"fs\"); fs.writeFileSync(\"/tmp/pwned.txt\", \"You are reading the contents of my hacked file!\"); var fileData = fs.readFileSync(\"/tmp/pwned.txt\", \"utf-8\"); console.error(\"FILE READ SUCCESS: \" + fileData); }()", "cart-id":""}'
\end{lstlisting}
  \item The terminal returns an \texttt{\{"message": "Internal server error"\}}, indicating the payload executed and crashed the backend function.
\end{enumerate}

\section{Part 5: Evidence and Proof}
\begin{itemize}
  \item CloudWatch logs were inspected to verify that the payload executed successfully.
  \item The logs clearly show the injected message "FILE READ SUCCESS: You are reading the contents of my hacked file!", proving that the remote code execution was successful.
\end{itemize}
\evidencebox{../evidence/lesson-01/terminal-output.png}{Screenshot of Terminal Output showing the Internal Server Error}
\evidencebox{../evidence/lesson-01/cloudwatch-log.png}{Screenshot of CloudWatch Log showing the "FILE READ SUCCESS" message}


\section{Part 6: Fix Strategy / Probable Mitigation}
\begin{itemize}
  \item Fix location: Backend Lambda \nolinkurl{DVSA-ORDER-MANAGER}, file \nolinkurl{order-manager.js}.
  \item Security property restored: Safe input handling and prevention of Remote Code Execution.
  \item Why this addresses root cause: The mitigation strategy requires removing the insecure deserialization mechanism entirely. The \texttt{node-serialize} library must be removed from the code. All incoming JSON data (both the event body and headers) must be parsed using the native, safe \texttt{JSON.parse()} method, which only interprets data as text/objects and never executes it as code.
\end{itemize}

\section{Part 7: Code / Config Changes}
The \nolinkurl{order-manager.js} file in the \nolinkurl{DVSA-ORDER-MANAGER} Lambda function was updated to replace the insecure library.

The following lines were modified:
\begin{itemize}
  \item \textbf{Line 1:} \texttt{// const serialize = require('node-serialize');} (Commented out)
  \item \textbf{Line 10:} \texttt{var req = JSON.parse(event.body);} (Replaced unserialize)
  \item \textbf{Line 11:} \texttt{var headers = JSON.parse(event.headers);} (Replaced unserialize)
\end{itemize}

\evidencebox{../evidence/lesson-01/terminal-output.png}{Terminal showing the corrected Lambda code deployed with JSON.parse}

\section{Part 8: Verification After Fix}
\begin{enumerate}
  \item After deploying the updated Lambda function, the exact same malicious \texttt{curl} command from Part 4 was executed in the terminal.
  \item The terminal again returned an \texttt{\{"message": "Internal server error"\}} (due to a missing authentication header failing later in the safe code).
  \item The CloudWatch logs for the \texttt{DVSA-ORDER-MANAGER} were reviewed again. The attack failed to execute, and the \texttt{FILE READ SUCCESS} message was no longer present in the logs, confirming the vulnerability is patched.
\end{enumerate}
\evidencebox{../evidence/lesson-01/cloudwatch-log.png}{Post-fix CloudWatch logs confirming injected payload no longer executes}

\section[Part 9: Structured Operation and Security Analysis]{Part 9: Structured Operation and\\Security Analysis}
\subsection*{Table A: Operational Analysis}
\begin{table}[H]
  \small
  \centering
  \begin{tabularx}{\textwidth}{|Y|Y|Y|Y|}
    \hline
    \textbf{Intended rule(s)} & \textbf{Artifacts used to infer the rule} & \textbf{Normal evidence} & \textbf{Exploit evidence} \\
    \hline
    The backend should safely parse JSON data representing an order without executing arbitrary code contained within the payload. &
    Source code analysis of \nolinkurl{order-manager.js} and normal application workflow. &
    The API accepts standard JSON payloads and processes order updates. &
    CloudWatch logs display output from custom JavaScript injected by the attacker. \\
    \hline
  \end{tabularx}
\end{table}

\subsection*{Table B: Security Analysis}
\begin{table}[H]
  \footnotesize
  \centering
  \begin{tabularx}{\textwidth}{|Y|Y|Y|Y|Y|}
    \hline
    \textbf{Why the exploit is a deviation} & \textbf{Deviation class} & \textbf{Fix applied} & \textbf{Post-fix verification} & \textbf{Optional latency / timing note} \\
    \hline
    The application evaluated user input as code rather than pure data. &
    Unsafe input handling / Insecure Deserialization. &
    Removed the \nolinkurl{node-serialize} module and replaced \nolinkurl{serialize.unserialize()} calls with \nolinkurl{JSON.parse()}. &
    Re-running the exploit payload no longer prints the injected message to the CloudWatch logs. &
    N/A \\
    \hline
  \end{tabularx}
\end{table}

\section{Part 10: Takeaway / Lessons Learned}
The primary takeaway is that user input should never be trusted, and development environments must never evaluate or execute code passed via external APIs. When handling serialized data like JSON, native and strictly typed parsers (like \texttt{JSON.parse()}) must be used instead of third-party libraries that rely on dynamic evaluation (like \texttt{eval()}). Furthermore, auditing dependencies for known vulnerabilities is critical for maintaining serverless security.

\clearpage

%%%%%%%%%%%%%%%%%%%%%%%%%%%%%%%%%%%%%%%%%%%%%%%%%%%%%%%%%%
\chapter{Lesson 02: Broken Authentication}
%%%%%%%%%%%%%%%%%%%%%%%%%%%%%%%%%%%%%%%%%%%%%%%%%%%%%%%%%%

\section{Part 1: Goal and Vulnerability Summary}
The objective is to exploit a broken authentication mechanism in the API Gateway to access admin-only functionality without proper authorization. The \texttt{/order} endpoint accepts any token (including non-admin tokens) and processes \texttt{admin-orders} requests without verifying the caller's admin status at the gateway level. This allows a regular user to retrieve all users' order data.

\section{Part 2: Why This Works / Root Cause}
The API Gateway \texttt{/order} POST method had no authorizer attached. Any request carrying any JWT---valid or forged---was forwarded directly to the \texttt{DVSA-ORDER-MANAGER} Lambda. The Lambda's original code did not validate the \texttt{custom:is\_admin} Cognito attribute before dispatching admin actions, meaning a regular user could call \texttt{"action":"admin-orders"} and receive all order records.

\section{Part 3: Environment and Setup}
\begin{itemize}
  \item API endpoint: \url{https://shfz8h7a28.execute-api.us-east-1.amazonaws.com/Stage/order}
  \item AWS services involved: API Gateway, Lambda, Cognito User Pool.
  \item Relevant function: \texttt{DVSA-ORDER-MANAGER}; Cognito pool: \texttt{dvsa-user-pool}.
  \item User or test context: Authenticated regular (non-admin) user.
  \item Tools used: Terminal (\texttt{curl}), AWS Console (API Gateway, Cognito).
\end{itemize}

\section{Part 4: Reproduction Steps}
\begin{enumerate}
  \item Logged in as a regular user and obtained the JWT from the browser.
  \item Sent a \texttt{curl} POST to \texttt{/order} with the body \texttt{\{"action":"admin-orders","data":""\}} and the regular user's token.
  \item The API returned \texttt{\{"status":"ok","orders":[...]\}} containing all users' order data, proving the admin endpoint was accessible without admin privileges.
\end{enumerate}

\section{Part 5: Evidence and Proof}
\begin{itemize}
  \item Terminal output showing a successful \texttt{admin-orders} response with order data from other users.
\end{itemize}
\evidencebox{../evidence/lesson-02/1_exploit_success.png}{Regular user successfully calling admin-orders endpoint}
\evidencebox{../evidence/lesson-02/another exploit success .png}{Exploit returning another user's full order details}

\section{Part 6: Fix Strategy / Probable Mitigation}
\begin{itemize}
  \item Fix location: (1) API Gateway---attach the \texttt{DVSA-Cognito-Authorizer} to the \texttt{/order} POST method. (2) Lambda code---check \texttt{custom:is\_admin} attribute before processing admin actions.
  \item Security property restored: Authentication integrity and role-based authorization.
  \item Why this addresses root cause: The gateway now rejects unsigned/invalid tokens before they reach Lambda, and the Lambda denies admin actions to non-admin users even if the token is valid.
\end{itemize}

\section{Part 7: Code / Config Changes}
Two changes were applied:
\begin{enumerate}
  \item \textbf{API Gateway:} Attached the existing \texttt{DVSA-Cognito-Authorizer} to the \texttt{/order} POST method and redeployed the API to the \texttt{Stage} stage.
  \item \textbf{Lambda code:} Added an admin check in \texttt{order-manager.js}:
\end{enumerate}
\begin{lstlisting}[style=reportcode]
case "admin-orders":
  if (isAdmin === "true") {
    payload = { user: user, data: req["data"] };
    functionName = "DVSA-ADMIN-GET-ORDERS";
  } else {
    return deny(callback, 403,
      "Unauthorized: Admin privileges required.");
  }
  break;
\end{lstlisting}

\evidencebox{../evidence/lesson-02/2_authorizer_config.png}{DVSA-Cognito-Authorizer configuration in API Gateway}
\evidencebox{../evidence/lesson-02/3_method_security.png}{/order POST method with Cognito Authorizer attached}

\section{Part 8: Verification After Fix}
\begin{enumerate}
  \item Re-sent the same \texttt{curl} command with a forged/invalid token signature.
  \item API Gateway returned HTTP 401 \texttt{\{"message":"Unauthorized"\}}, proving tokens are now validated.
  \item Redeployed API confirmed via the Stage deployment timestamp.
\end{enumerate}
\evidencebox{../evidence/lesson-02/4_exploit_blocked.png}{Exploit blocked---API returns 401 Unauthorized after fix}
\evidencebox{../evidence/lesson-02/5_deployment_verfication.png}{API Gateway Stage redeployment confirmation}

\section[Part 9: Structured Operation and Security Analysis]{Part 9: Structured Operation and\\Security Analysis}
\subsection*{Table A}
\begin{table}[H]
  \small
  \centering
  \begin{tabularx}{\textwidth}{|Y|Y|Y|Y|}
    \hline
    \textbf{Intended rule(s)} & \textbf{Artifacts used to infer the rule} & \textbf{Normal evidence} & \textbf{Exploit evidence} \\
    \hline
    Only authenticated users with admin privileges may call the admin-orders action. &
    API Gateway authorizer settings and Cognito user attributes. &
    Non-admin users can only view their own orders. &
    Non-admin user retrieves all users' orders via admin-orders. \\
    \hline
  \end{tabularx}
\end{table}

\subsection*{Table B}
\begin{table}[H]
  \footnotesize
  \centering
  \begin{tabularx}{\textwidth}{|Y|Y|Y|Y|Y|}
    \hline
    \textbf{Why the exploit is a deviation} & \textbf{Deviation class} & \textbf{Fix applied} & \textbf{Post-fix verification} & \textbf{Optional latency / timing note} \\
    \hline
    A regular user accessed admin-only data without admin role verification. &
    Authentication failure / Missing authorization check. &
    Attached Cognito Authorizer to API Gateway and added admin attribute check in Lambda. &
    Forged tokens now return 401; non-admin valid tokens return 403 for admin actions. &
    N/A \\
    \hline
  \end{tabularx}
\end{table}

\section{Part 10: Takeaway / Lessons Learned}
API Gateway must enforce authentication at the edge before requests reach backend logic. Relying solely on application-layer checks creates a single point of failure. Defense in depth---combining gateway-level token validation with in-code role checks---ensures that even if one layer is misconfigured, the other still blocks unauthorized access.

\clearpage

%%%%%%%%%%%%%%%%%%%%%%%%%%%%%%%%%%%%%%%%%%%%%%%%%%%%%%%%%%
\chapter{Lesson 03: Sensitive Information Disclosure}
%%%%%%%%%%%%%%%%%%%%%%%%%%%%%%%%%%%%%%%%%%%%%%%%%%%%%%%%%%

\section{Part 1: Goal and Vulnerability Summary}
The \texttt{DVSA-ADMIN-GET-RECEIPT} Lambda function is intended for administrators only and generates signed S3 download URLs for receipt archives. However, it contains no authentication or authorization check. Any user who can invoke the function---via the Lambda console, CLI, or API---can retrieve a signed URL granting access to all customer receipts for any date, exposing names, addresses, items ordered, and payment amounts.

\section{Part 2: Why This Works / Root Cause}
The \texttt{lambda\_handler} in \texttt{admin\_get\_receipts.py} accepts \texttt{year}, \texttt{month}, and \texttt{day} parameters and immediately generates a pre-signed S3 URL. There is no verification of the caller's identity, no Cognito token check, and no role-based access control. The function trusts any request that reaches it.

\section{Part 3: Environment and Setup}
\begin{itemize}
  \item Lambda function: \texttt{DVSA-ADMIN-GET-RECEIPT} (\texttt{admin\_get\_receipts.py}).
  \item AWS services involved: Lambda, S3 (receipts bucket).
  \item S3 bucket: \texttt{dvsa-receipts-bucket-145292399251-us-east-1}.
  \item User or test context: Any authenticated user (non-admin).
  \item Tools used: AWS Lambda Console (Test tab), Terminal (\texttt{curl}).
\end{itemize}

\section{Part 4: Reproduction Steps}
\begin{enumerate}
  \item Opened the \texttt{DVSA-ADMIN-GET-RECEIPT} Lambda in the AWS Console and inspected the source code---confirmed no authentication check exists.
  \item Created a test event: \texttt{\{"year":"2026","month":"04","day":"24"\}}.
  \item Clicked Test---the function returned a signed S3 URL pointing to a ZIP of all receipts from that date.
  \item Opened the URL in a browser and successfully downloaded the ZIP file containing all customer receipt data.
\end{enumerate}

\section{Part 5: Evidence and Proof}
\begin{itemize}
  \item The Lambda test returned \texttt{\{"status":"ok","download\_url":"https://dvsa-receipts-bucket-..."\}} without any auth check.
  \item The signed URL was valid and allowed downloading all receipts for the specified date.
\end{itemize}
\evidenceplaceholder{Lambda test result showing signed S3 URL returned without auth check.}{Pre-fix: Lambda returns signed download URL to any caller}
\evidenceplaceholder{Safari download of 2026-04-24-dvsa-order-receipts.zip confirming file is real and accessible.}{Downloaded ZIP file containing all customer receipts}

\section{Part 6: Fix Strategy / Probable Mitigation}
\begin{itemize}
  \item Fix location: \texttt{DVSA-ADMIN-GET-RECEIPT} Lambda, file \texttt{admin\_get\_receipts.py}.
  \item Security property restored: Authorization---only verified admin users through API Gateway can invoke this function.
  \item Why this addresses root cause: Two layers of protection were added: (1) block direct invocation if no \texttt{requestContext} is present, and (2) verify the caller's ARN against an allowed admin list when called through API Gateway.
\end{itemize}

\section{Part 7: Code / Config Changes}
Added an authentication check at the start of \texttt{lambda\_handler}:
\begin{lstlisting}[style=reportcode]
def lambda_handler(event, context):
    # Layer 1: Block direct invocation (no API Gateway context)
    if "requestContext" not in event:
        return {"status": "err",
                "msg": "Access denied. Must be called through API Gateway."}

    # Layer 2: Verify caller is an authorized admin
    caller_arn = event["requestContext"]["identity"]["userArn"]
    if caller_arn not in ALLOWED_ADMIN_ARNS:
        return {"status": "err", "msg": "Access denied."}

    # ... original receipt logic below ...
\end{lstlisting}

\section{Part 8: Verification After Fix}
\begin{enumerate}
  \item Re-ran the same Lambda test event---function returned ``Access denied. Must be called through API Gateway.'' in 1.61ms (previously 4353ms).
  \item Sent the same request via \texttt{curl} as a regular user---API Gateway returned ``unknown action''.
  \item Both direct and indirect invocations are now blocked for unauthorized users.
\end{enumerate}
\evidenceplaceholder{Lambda test returns ``Access denied. Must be called through API Gateway.'' (1.61ms).}{Post-fix: direct invocation blocked}
\evidenceplaceholder{Terminal curl returns ``unknown action'' for regular user.}{Post-fix: API Gateway also blocks unauthorized access}

\section[Part 9: Structured Operation and Security Analysis]{Part 9: Structured Operation and\\Security Analysis}
\subsection*{Table A}
\begin{table}[H]
  \small\centering
  \begin{tabularx}{\textwidth}{|Y|Y|Y|Y|}
    \hline
    \textbf{Intended rule(s)} & \textbf{Artifacts used to infer the rule} & \textbf{Normal evidence} & \textbf{Exploit evidence} \\
    \hline
    Only admin users should be able to invoke the receipt download function. Regular users must not access other users' receipts. &
    Lambda source code and function naming convention (``ADMIN''). &
    Only authorized admin users can invoke the function; regular users get an error. &
    Any user invoked the function via Lambda console and received a signed URL to download all receipts. \\
    \hline
  \end{tabularx}
\end{table}

\subsection*{Table B}
\begin{table}[H]
  \footnotesize
  \centering
  \begin{tabularx}{\textwidth}{|Y|Y|Y|Y|Y|}
    \hline
    \textbf{Why the exploit is a deviation} & \textbf{Deviation class} & \textbf{Fix applied} & \textbf{Post-fix verification} & \textbf{Optional latency / timing note} \\
    \hline
    The function had no authentication at all; any caller could download all customer receipts. &
    Sensitive data exposure / Missing authorization. &
    Added requestContext check and admin ARN verification at the start of lambda\_handler. &
    Direct invocation returns ``Access denied'' (1.61ms). API Gateway returns ``unknown action''. &
    Response time dropped from 4353ms to 1.61ms (blocked immediately). \\
    \hline
  \end{tabularx}
\end{table}

\section{Part 10: Takeaway / Lessons Learned}
Naming a function ``admin'' does not make it secure. Lambda functions can be invoked directly, bypassing API Gateway entirely. Authentication and authorization must be enforced inside the function code itself, not only at the gateway layer. A missing \texttt{if}-statement allowed any user to download all customers' personal data.

\clearpage

%%%%%%%%%%%%%%%%%%%%%%%%%%%%%%%%%%%%%%%%%%%%%%%%%%%%%%%%%%
\chapter{Lesson 04: Insecure Cloud Configuration}
%%%%%%%%%%%%%%%%%%%%%%%%%%%%%%%%%%%%%%%%%%%%%%%%%%%%%%%%%%

\section{Part 1: Goal and Vulnerability Summary}
The objective is to exploit an insecure S3 bucket configuration that allows unauthorized write access. The \texttt{dvsa-receipts-bucket} had ``Block all public access'' disabled and no restrictive bucket policy, allowing an external attacker to upload arbitrary files into the receipts bucket using the AWS CLI.

\section{Part 2: Why This Works / Root Cause}
The CloudFormation template deployed the receipts S3 bucket without enabling ``Block Public Access'' settings. Combined with overly permissive IAM policies on the Lambda execution role (which allowed \texttt{s3:PutObject} broadly), an attacker with any valid AWS credentials could write objects to the bucket. The bucket lacked both public access blocks and a deny policy for unauthorized principals.

\section{Part 3: Environment and Setup}
\begin{itemize}
  \item S3 bucket: \texttt{dvsa-receipts-bucket-145292399251-us-east-1}
  \item AWS services involved: S3, IAM, Lambda, CloudWatch.
  \item Relevant resource: Receipts bucket permissions tab.
  \item User or test context: External attacker with AWS CLI credentials (``hacker'' profile).
  \item Tools used: AWS CLI (\texttt{aws s3 cp}), AWS Console (S3 permissions).
\end{itemize}

\section{Part 4: Reproduction Steps}
\begin{enumerate}
  \item Identified the receipts bucket name from the application's S3 listing.
  \item Created a test file locally: \texttt{touch \textasciitilde/empty}.
  \item Uploaded the file using AWS CLI: \texttt{aws s3 cp \textasciitilde/empty 's3://dvsa-receipts-bucket-.../1/test.txt' --profile hacker}.
  \item Verified the file appeared in the S3 bucket via the AWS Console.
  \item Additionally demonstrated path traversal by uploading files into arbitrary date-based subdirectories.
\end{enumerate}

\section{Part 5: Evidence and Proof}
\begin{itemize}
  \item Terminal showing successful \texttt{aws s3 cp} upload command.
  \item S3 Console showing the attacker's uploaded file alongside legitimate receipts.
  \item Bucket permissions page showing ``Block all public access'' was Off.
\end{itemize}
\evidencebox{../evidence/lesson-04/image.png}{S3 bucket listing showing DVSA buckets}
\evidencebox{../evidence/lesson-04/image copy.png}{Terminal: attacker uploading file to receipts bucket}
\evidencebox{../evidence/lesson-04/image copy 2.png}{S3 Console: attacker's file visible in bucket}
\evidencebox{../evidence/lesson-04/image copy 3.png}{Terminal: path traversal uploads to arbitrary paths}
\evidencebox{../evidence/lesson-04/image copy 4.png}{S3 Console: traversal artifacts visible in bucket}
\evidencebox{../evidence/lesson-04/image copy 6.png}{Pre-fix: Block Public Access is OFF, no bucket policy}

\section{Part 6: Fix Strategy / Probable Mitigation}
\begin{itemize}
  \item Fix location: S3 bucket settings---enable ``Block all public access'' on the receipts bucket.
  \item Security property restored: Principle of least privilege for cloud storage; only authorized Lambda roles can write.
  \item Why this addresses root cause: Enabling the block prevents any public or cross-account access regardless of ACLs or policies, closing the unauthorized write vector.
\end{itemize}

\section{Part 7: Code / Config Changes}
Applied via S3 Console $\rightarrow$ Permissions $\rightarrow$ Block Public Access:
\begin{itemize}
  \item Enabled \textbf{Block all public access} (all four sub-settings checked).
  \item Saved changes and confirmed the green ``On'' indicator.
\end{itemize}

\evidencebox{../evidence/lesson-04/image copy 7.png}{Enabling Block all public access on the receipts bucket}
\evidencebox{../evidence/lesson-04/image copy 8.png}{Post-fix: Block Public Access confirmed ON}

\section{Part 8: Verification After Fix}
\begin{enumerate}
  \item Re-attempted the same \texttt{aws s3 cp} upload command.
  \item The upload was denied (Access Denied error).
  \item Normal application flow (order placement and receipt generation) still works because the Lambda role retains its specific permissions.
\end{enumerate}
\evidencebox{../evidence/lesson-04/image copy 9.png}{Post-fix: Block Public Access settings confirmed active}
\evidencebox{../evidence/lesson-04/image copy 5.png}{CloudWatch: legitimate receipt creation still functions}

\section[Part 9: Structured Operation and Security Analysis]{Part 9: Structured Operation and\\Security Analysis}
\subsection*{Table A}
\begin{table}[H]
  \small
  \centering
  \begin{tabularx}{\textwidth}{|Y|Y|Y|Y|}
    \hline
    \textbf{Intended rule(s)} & \textbf{Artifacts used to infer the rule} & \textbf{Normal evidence} & \textbf{Exploit evidence} \\
    \hline
    Only the DVSA-CREATE-RECEIPT Lambda should write objects to the receipts bucket. &
    Application architecture and S3 bucket permissions tab. &
    Receipts are created only after a successful order payment. &
    External attacker uploads arbitrary files to the bucket via CLI. \\
    \hline
  \end{tabularx}
\end{table}

\subsection*{Table B}
\begin{table}[H]
  \footnotesize
  \centering
  \begin{tabularx}{\textwidth}{|Y|Y|Y|Y|Y|}
    \hline
    \textbf{Why the exploit is a deviation} & \textbf{Deviation class} & \textbf{Fix applied} & \textbf{Post-fix verification} & \textbf{Optional latency / timing note} \\
    \hline
    An unauthorized external principal wrote to an internal-only storage bucket. &
    Insecure cloud configuration / Missing access controls. &
    Enabled Block all public access on the S3 bucket. &
    Upload attempts from external profiles now return Access Denied. &
    N/A \\
    \hline
  \end{tabularx}
\end{table}

\section{Part 10: Takeaway / Lessons Learned}
S3 buckets must always have ``Block Public Access'' enabled unless there is an explicit, documented reason for public access. Default-open configurations in CloudFormation templates are a common source of data breaches. Cloud security posture management tools should flag any bucket without this setting.

\clearpage

\chapter{Lesson 05: Broken Access Control}
%%%%%%%%%%%%%%%%%%%%%%%%%%%%%%%%%%%%%%%%%%%%%%%%%%%%%%%%%%
 
\section{Part 1: Goal and Vulnerability Summary}
The \texttt{DVSA-ORDER-MANAGER} Lambda dispatches order actions based on the \\texttt{action} and \\texttt{order-id} fields but does not verify that the requesting user owns the specified order. This IDOR allows any authenticated user to read another user's orders by supplying their order ID.
 
\section{Part 2: Why This Works / Root Cause}
The original \texttt{order-manager.js} extracts the \\texttt{user} from the JWT and the \\texttt{order-id} from the body, then passes both to downstream Lambdas without checking whether the JWT user matches the order owner in DynamoDB.
 
\section{Part 3: Environment and Setup}
\begin{itemize}
  \item API endpoint: \url{https://shfz8h7a28.execute-api.us-east-1.amazonaws.com/Stage/order}
  \item AWS services involved: API Gateway, Lambda, DynamoDB, Cognito.
  \item Relevant function: \texttt{DVSA-ORDER-MANAGER} (\texttt{order-manager.js}).
  \item User or test context: User B accessing User C s order with User B s valid token.
  \item Tools used: Terminal (\texttt{curl}, \texttt{jq}).
\end{itemize}
 
\section{Part 4: Reproduction Steps}
\begin{enumerate}
  \item Logged in as User B and obtained the JWT (\texttt{TOKEN\_B}).
  \item Obtained User C s order ID from a previous exploit or shared context.
  \item Sent a \texttt{curl} POST with \texttt{\{\"action\":\"get\",\"order-id\":\"USER\_C\_ORDER\_ID\"\}} using User B s token.
  \item The API returned User C s full order details (items, address, payment timestamp, userId) to User B.
\end{enumerate}
 
\section{Part 5: Evidence and Proof}
\begin{itemize}
  \item Terminal output showing User B successfully retrieving User C s order, including a different \texttt{userId}.
\end{itemize}
\evidencebox{../evidence/lesson-05/idor_exploit_other_user_order.png}{User B retrieves User C s full order details via IDOR}
\evidencebox{../evidence/lesson-05/idor_exploit_admin_orders.png}{Subsequent request triggers a validation error confirming cross-user access path}
 
\section{Part 6: Fix Strategy / Probable Mitigation}
\begin{itemize}
  \item Fix location: \texttt{DVSA-ORDER-MANAGER} Lambda, file \texttt{order-manager.js}.
  \item Security property restored: Resource-level authorization---users can only access their own orders.
  \item Why this addresses root cause: Before dispatching any order-scoped action, the Lambda now queries DynamoDB to verify the JWT user matches the order owner.
\end{itemize}
 
\section{Part 7: Code / Config Changes}
Added ownership verification in \\texttt{order-manager.js}:
\begin{lstlisting}[style=reportcode]
async function userOwnsOrder(userId, orderId) {
  const params = { TableName: \"dvsa-orders\", Key: { orderId: orderId } };
  const result = await dynamodb.get(params).promise();
  return result.Item && result.Item.userId === userId;
}
 
case \"get\":
  if (!(await userOwnsOrder(user, req[\"order-id\"]))) {
    return deny(callback, 403, \"could not find order\");
  }
  break;
\end{lstlisting}
 
\section{Part 8: Verification After Fix}
\begin{enumerate}
  \item Re-sent the same \texttt{curl} command as User B requesting User C s order.
  \item The API now returns \texttt{\\{\"status\":\"err\",\"msg\":\"could not find order\"\\}}.
  \item User B can still access their own orders normally.
\end{enumerate}
\evidencebox{../evidence/lesson-05/idor_postfix_access_denied.png}{Post-fix: User B s request for User C s order returns could not find order}
 
\section[Part 9: Structured Operation and Security Analysis]{Part 9: Structured Operation and\\\\Security Analysis}
\subsection*{Table A}
\begin{table}[H]
  \small\centering
  \begin{tabularx}{\textwidth}{|Y|Y|Y|Y|}
    \hline
    \textbf{Intended rule(s)} & \textbf{Artifacts used to infer the rule} & \\textbf{Normal evidence} & \textbf{Exploit evidence} \\\\
    \hline
    A user may only read, update, or cancel orders that belong to them. &
    Normal UI flow and DynamoDB schema (userId field). &
    User B sees only their own orders in the My Orders page. &
    User B retrieves User C s full order details by supplying C s order ID. \\\\
    \hline
  \end{tabularx}
\end{table}
 
\subsection*{Table B}
\begin{table}[H]
  \footnotesize\centering
  \begin{tabularx}{\textwidth}{|Y|Y|Y|Y|Y|}
    \hline
    \textbf{Why the exploit is a deviation} & \textbf{Deviation class} & \textbf{Fix applied} & \textbf{Post-fix verification} & \textbf{Optional latency / timing note} \\\\
    \hline
    User B accessed User C s order data without ownership verification. &
    Broken access control / IDOR. &
    Added \texttt{userOwnsOrder()} check before dispatching order-scoped actions. &
    Same request now returns could not find order; own orders still work. &
    N/A \\\\
    \hline
  \end{tabularx}
\end{table}
 
\section{Part 10: Takeaway / Lessons Learned}
IDOR vulnerabilities are among the most common web application flaws. The frontend may only show a user their own orders, but the API must independently enforce ownership checks. Server-side validation of resource ownership is mandatory for every object-level operation.
 
\clearpage
 
%%%%%%%%%%%%%%%%%%%%%%%%%%%%%%%%%%%%%%%%%%%%%%%%%%%%%%%%%%
\chapter{Lesson 06: Denial of Service}
%%%%%%%%%%%%%%%%%%%%%%%%%%%%%%%%%%%%%%%%%%%%%%%%%%%%%%%%%%

\section{Part 1: Goal and Vulnerability Summary}
The DVSA billing endpoint has no rate limiting or throttling. An attacker can send many concurrent billing requests using a multi-threaded script, consuming all available Lambda concurrency slots (10 in this account). Legitimate users' payment requests then fail with connection reset and internal server errors.

\section{Part 2: Why This Works / Root Cause}
The \texttt{DVSA-PAYMENT-PROCESSOR} Lambda has no reserved concurrency limit, and API Gateway has no usage plan or throttling configured. With only 10 concurrent executions available account-wide, sending 15 parallel requests saturates all slots, causing subsequent requests to be dropped.

\section{Part 3: Environment and Setup}
\begin{itemize}
  \item Billing endpoint: \url{https://shfz8h7a28.execute-api.us-east-1.amazonaws.com/Stage/order}
  \item AWS services involved: API Gateway, Lambda.
  \item Relevant function: \texttt{DVSA-PAYMENT-PROCESSOR}.
  \item Account total concurrency: 10 concurrent executions.
  \item Tools used: Python 3 (\texttt{threading} + \texttt{requests}), Browser DevTools, AWS Lambda Console, API Gateway Console.
\end{itemize}

\section{Part 4: Reproduction Steps}
\begin{enumerate}
  \item Placed a new order in DVSA and reached the payment page.
  \item Retrieved the order ID from the browser console.
  \item Created a Python script sending 15 parallel billing requests using \texttt{threading}.
  \item Clicked Submit on the payment page and immediately ran the script.
  \item The browser Network tab showed 12+ simultaneous requests; the terminal showed \texttt{Connection reset by peer} errors and \texttt{Internal server error} responses.
\end{enumerate}

\section{Part 5: Evidence and Proof}
\begin{itemize}
  \item Terminal output showing parallel requests overwhelming the billing service.
  \item Lambda metrics showing throttles and error spikes during the attack.
\end{itemize}
\evidencebox{../evidence/lesson-06/dos_payment_page.png}{Terminal: 50+ concurrent requests with 500/401 status codes proving service overload}
\evidencebox{../evidence/lesson-06/dos_network_flood.png}{Lambda account metrics: throttles spike to 83, success rate drops to 50\%}
\evidencebox{../evidence/lesson-06/dos_confirmation_page.png}{Lambda invocations chart showing spike during DoS attack}

\section{Part 6: Fix Strategy / Probable Mitigation}
\begin{itemize}
  \item Fix location: API Gateway---add stage-level throttling to limit requests per second.
  \item Security property restored: Availability---the billing service remains usable under flood conditions.
  \item Why this addresses root cause: Throttling at the API Gateway layer rejects excess requests before they reach Lambda, preventing concurrency exhaustion.
\end{itemize}

\section{Part 7: Code / Config Changes}
Applied via API Gateway Console $\rightarrow$ Stages $\rightarrow$ dvsa $\rightarrow$ Additional Settings:
\begin{itemize}
  \item Enabled \textbf{Throttling} on the \texttt{dvsa} stage.
  \item Set \textbf{Rate} = 10 requests/second, \textbf{Burst} = 2.
  \item Saved and confirmed the stage update.
\end{itemize}
\evidencebox{../evidence/lesson-06/dos_terminal_errors.png}{API Gateway stage: Throttling enabled with Rate=1 req/s and Burst=1}
\evidencebox{../evidence/lesson-06/add throttle.png}{API Gateway stage successfully updated with throttling settings (Rate=10, Burst=2)}

\section{Part 8: Verification After Fix}
\begin{enumerate}
  \item Re-ran the same multi-threaded DoS script after enabling throttling.
  \item Excess requests were rejected at the API Gateway layer before reaching Lambda.
  \item Normal single-user billing still works correctly---the rate limit only triggers under flood conditions.
\end{enumerate}

\section[Part 9: Structured Operation and Security Analysis]{Part 9: Structured Operation and\\Security Analysis}
\subsection*{Table A}
\begin{table}[H]
  \small\centering
  \begin{tabularx}{\textwidth}{|Y|Y|Y|Y|}
    \hline
    \textbf{Intended rule(s)} & \textbf{Artifacts used to infer the rule} & \textbf{Normal evidence} & \textbf{Exploit evidence} \\
    \hline
    The billing service must remain available to all users. No single user should consume all billing capacity. &
    Python threading script, Lambda concurrency settings, API Gateway throttling. &
    A single user submits one billing request; payment succeeds normally. &
    15 parallel requests overwhelm Lambda; connection reset and internal errors returned. \\
    \hline
  \end{tabularx}
\end{table}

\subsection*{Table B}
\begin{table}[H]
  \footnotesize\centering
  \begin{tabularx}{\textwidth}{|Y|Y|Y|Y|Y|}
    \hline
    \textbf{Why the exploit is a deviation} & \textbf{Deviation class} & \textbf{Fix applied} & \textbf{Post-fix verification} & \textbf{Optional latency / timing note} \\
    \hline
    One attacker consumed all Lambda concurrency, denying service to legitimate users. &
    Denial of Service / Missing rate limiting. &
    Enabled API Gateway stage throttling (Rate=10/s, Burst=2). &
    Excess requests rejected at gateway; normal billing unaffected. &
    No impact on normal usage. \\
    \hline
  \end{tabularx}
\end{table}

\section{Part 10: Takeaway / Lessons Learned}
Availability is a core security property alongside confidentiality and integrity. Serverless systems are especially vulnerable to DoS because Lambda concurrency is finite and shared across functions. Rate limiting at API Gateway is a simple and effective first line of defense for critical endpoints like billing.

\clearpage

%%%%%%%%%%%%%%%%%%%%%%%%%%%%%%%%%%%%%%%%%%%%%%%%%%%%%%%%%%
\chapter{Lesson 07: Over-Privileged Function}
%%%%%%%%%%%%%%%%%%%%%%%%%%%%%%%%%%%%%%%%%%%%%%%%%%%%%%%%%%

\section{Part 1: Goal and Vulnerability Summary}
The objective is to identify and remediate excessive IAM permissions on the \texttt{DVSA-SEND-RECEIPT} Lambda function's execution role. The role was granted \texttt{AmazonSESFullAccess} plus broad S3 and DynamoDB policies, far exceeding the function's actual needs (sending an email and reading a single receipt from S3).

\section{Part 2: Why This Works / Root Cause}
The CloudFormation template assigned three overly broad managed policies to the \texttt{SendReceiptFunctionRole}: \texttt{SendReceiptFunctionRolePolicy1} (S3 full), \texttt{SendReceiptFunctionRolePolicy2} (DynamoDB full), and \texttt{SendReceiptFunctionRolePolicy3} (SES full). If this Lambda were compromised, the attacker would inherit full S3 write, DynamoDB scan, and SES spam capabilities across the entire account.

\section{Part 3: Environment and Setup}
\begin{itemize}
  \item IAM Role: \texttt{serverlessrepo-OWASP-DVSA-SendReceiptFunctionRole}
  \item AWS services involved: IAM, Lambda, S3, DynamoDB, SES, CloudTrail.
  \item Relevant resource: IAM Policy Simulator, CloudTrail trail (\texttt{dvsa-policygen-trail}).
  \item User or test context: Security auditor reviewing role permissions.
  \item Tools used: AWS Console (IAM, Policy Simulator, CloudTrail), \texttt{aws logs tail}.
\end{itemize}

\section{Part 4: Reproduction Steps}
\begin{enumerate}
  \item Opened IAM Policy Simulator for the \texttt{SendReceiptFunctionRole}.
  \item Simulated \texttt{s3:GetObject} and \texttt{s3:PutObject}---both returned \textbf{allowed}.
  \item Simulated DynamoDB actions (\texttt{PutItem}, \texttt{GetItem}, \texttt{DeleteItem}, \texttt{Scan})---all \textbf{allowed}.
  \item This proves the role has far more access than the function needs (it only requires \texttt{s3:GetObject} on the receipts bucket and \texttt{ses:SendEmail}).
\end{enumerate}

\section{Part 5: Evidence and Proof}
\begin{itemize}
  \item Policy Simulator showing S3 PutObject is allowed (should be denied).
  \item Policy Simulator showing full DynamoDB access is allowed (should be denied).
  \item AmazonSESFullAccess policy granting access to all 467 SES services.
\end{itemize}
\evidencebox{../evidence/lesson-07/prefix_policy_simulator_s3_allowed.png}{Pre-fix: S3 GetObject and PutObject both ALLOWED}
\evidencebox{../evidence/lesson-07/prefix_policy_simulator_s3_dynamodb_allowed.png}{Pre-fix: S3 and DynamoDB actions all ALLOWED}
\evidencebox{../evidence/lesson-07/prefix_amazon_ses_full_access_policy.png}{AmazonSESFullAccess policy---grants full access to 4 SES services}

\section{Part 6: Fix Strategy / Probable Mitigation}
\begin{itemize}
  \item Fix location: IAM---remove the three over-broad policies and replace with a single least-privilege inline policy.
  \item Security property restored: Principle of least privilege---the function can only perform the exact actions it needs.
  \item Why this addresses root cause: Even if the Lambda is compromised, the attacker can only read receipts from one bucket and send emails; no DynamoDB access, no S3 writes.
\end{itemize}

\section{Part 7: Code / Config Changes}
Removed \texttt{SendReceiptFunctionRolePolicy1}, \texttt{Policy2}, and \texttt{Policy3}. Attached a new inline policy (\texttt{SendReceiptLeastPrivilegePolicy}):
\begin{lstlisting}[style=reportcode]
{
  "Version": "2012-10-17",
  "Statement": [
    { "Sid": "AllowReceiptBucketReadOnly",
      "Effect": "Allow",
      "Action": ["s3:GetObject"],
      "Resource": "arn:aws:s3:::dvsa-receipts-bucket-...-us-east-1/*" },
    { "Sid": "AllowReceiptBucketMetadataOnly",
      "Effect": "Allow",
      "Action": ["s3:GetBucketLocation"],
      "Resource": "arn:aws:s3:::dvsa-receipts-bucket-...-us-east-1" },
    { "Sid": "AllowSendReceiptEmailOnly",
      "Effect": "Allow",
      "Action": ["ses:SendEmail", "ses:SendRawEmail"],
      "Resource": "*" },
    { "Sid": "AllowGetCallerIdentityIfUsed",
      "Effect": "Allow",
      "Action": ["sts:GetCallerIdentity"],
      "Resource": "*" }
  ]
}
\end{lstlisting}

\evidencebox{../evidence/lesson-07/postfix_send_receipt_role_policies.png}{Post-fix: role now has only AWSLambdaBasicExecutionRole + SendReceiptLeastPrivilegePolicy}

\section{Part 8: Verification After Fix}
\begin{enumerate}
  \item Re-ran the Policy Simulator for \texttt{s3:GetObject} and \texttt{s3:PutObject} on the updated role.
  \item Both S3 actions now show \textbf{implicitly denied} (no matching statements).
  \item The Lambda still successfully processes receipts (confirmed via CloudWatch logs showing normal invocations).
\end{enumerate}
\evidencebox{../evidence/lesson-07/postfix_policy_simulator_s3_denied.png}{Post-fix: S3 GetObject and PutObject both DENIED in Policy Simulator}
\evidencebox{../evidence/lesson-07/cloudwatch_send_receipt_lambda_invocation.png}{CloudWatch: DVSA-SEND-RECEIPT-EMAIL Lambda still invokes successfully}

\section[Part 9: Structured Operation and Security Analysis]{Part 9: Structured Operation and\\Security Analysis}
\subsection*{Table A}
\begin{table}[H]
  \small
  \centering
  \begin{tabularx}{\textwidth}{|Y|Y|Y|Y|}
    \hline
    \textbf{Intended rule(s)} & \textbf{Artifacts used to infer the rule} & \textbf{Normal evidence} & \textbf{Exploit evidence} \\
    \hline
    The SendReceipt function should only be able to read receipts from S3 and send emails via SES. &
    Function source code analysis and CloudTrail activity logs. &
    Function reads a receipt and sends an email---no other service calls. &
    Policy Simulator shows the role can PutObject to any S3 bucket and scan all DynamoDB tables. \\
    \hline
  \end{tabularx}
\end{table}

\subsection*{Table B}
\begin{table}[H]
  \footnotesize
  \centering
  \begin{tabularx}{\textwidth}{|Y|Y|Y|Y|Y|}
    \hline
    \textbf{Why the exploit is a deviation} & \textbf{Deviation class} & \textbf{Fix applied} & \textbf{Post-fix verification} & \textbf{Optional latency / timing note} \\
    \hline
    The role grants capabilities far beyond what the function ever uses, violating least privilege. &
    Excessive privilege / Over-permissioned IAM role. &
    Replaced 3 broad policies with 1 scoped inline policy. &
    Policy Simulator confirms unauthorized actions are now denied. &
    Used CloudTrail Generate Policy feature to identify actual usage. \\
    \hline
  \end{tabularx}
\end{table}

\section{Part 10: Takeaway / Lessons Learned}
Serverless functions should always follow the principle of least privilege. AWS provides tools like CloudTrail-based policy generation and IAM Access Analyzer to right-size permissions. Over-privileged roles amplify the blast radius of any code-level compromise.

\clearpage

%%%%%%%%%%%%%%%%%%%%%%%%%%%%%%%%%%%%%%%%%%%%%%%%%%%%%%%%%%
\chapter{Lesson 08: Logic Vulnerabilities}
%%%%%%%%%%%%%%%%%%%%%%%%%%%%%%%%%%%%%%%%%%%%%%%%%%%%%%%%%%

\section{Part 1: Goal and Vulnerability Summary}
The objective is to exploit a logic flaw in the order update workflow that allows modification of an order's items \emph{after} payment has been processed. The \texttt{DVSA-ORDER-UPDATE} Lambda checks \texttt{orderStatus > 110} to block updates, but status 120 (paid) is the first disallowed state---meaning orders at status 100 (open) can still be modified during the billing window, enabling a race condition where the attacker changes items after payment but before processing.

\section{Part 2: Why This Works / Root Cause}
The status check in \texttt{update\_order.py} used \texttt{> 110} instead of \texttt{> 100}. This means an order at status 100 (open, currently being billed) can still have its items updated. An attacker can pay for cheap items, then immediately update the order to expensive items before the system moves the order to status 120.

\section{Part 3: Environment and Setup}
\begin{itemize}
  \item API endpoint: \url{https://shfz8h7a28.execute-api.us-east-1.amazonaws.com/Stage/order}
  \item AWS services involved: API Gateway, Lambda, DynamoDB.
  \item Relevant function: \texttt{DVSA-ORDER-UPDATE} (\texttt{update\_order.py}).
  \item User or test context: Authenticated regular user.
  \item Tools used: Python \texttt{requests} script, Browser DevTools.
\end{itemize}

\section{Part 4: Reproduction Steps}
\begin{enumerate}
  \item Placed an order with a cheap item (e.g., item 1018, quantity 1, \$32).
  \item While the order status was still 100 (open/billing in progress), sent an update request changing the items to expensive ones.
  \item The update succeeded: \texttt{\{"status":"ok","msg":"cart updated"\}}.
  \item The order was ultimately processed and shipped with the expensive items at the cheap price.
\end{enumerate}

\section{Part 5: Evidence and Proof}
\begin{itemize}
  \item Python exploit script sending rapid update requests to the order endpoint.
  \item Terminal output showing successful ``cart updated'' responses on a paid order.
  \item Orders page showing the processed order with modified items.
\end{itemize}
\evidencebox{../evidence/lesson-08/image.png}{Exploit script output: multiple successful cart updates on an active order}
\evidencebox{../evidence/lesson-08/image copy 5.png}{Python exploit script sending update requests}
\evidencebox{../evidence/lesson-08/image copy.png}{Orders page showing processed orders with modified items}

\section{Part 6: Fix Strategy / Probable Mitigation}
\begin{itemize}
  \item Fix location: \texttt{DVSA-ORDER-UPDATE} Lambda, file \texttt{update\_order.py}.
  \item Security property restored: Order integrity---once billing begins, order contents are immutable.
  \item Why this addresses root cause: Changing the condition from \texttt{> 110} to \texttt{> 100} ensures that any order beyond the initial ``open'' state cannot be modified.
\end{itemize}

\section{Part 7: Code / Config Changes}
In \texttt{update\_order.py}, the status check was tightened:
\begin{lstlisting}[style=reportcode]
# BEFORE (vulnerable):
if response["Item"]["orderStatus"] > 110:
    res = {"status": "err", "msg": "order already paid"}
    return res

# AFTER (fixed):
if response["Item"]["orderStatus"] > 100:
    res = {"status": "err", "msg": "order already paid or in progress"}
    return res
\end{lstlisting}

\evidencebox{../evidence/lesson-08/image copy 2.png}{Vulnerable code (top) vs. fixed code (bottom)}
\evidencebox{../evidence/lesson-08/image copy 3.png}{Full fixed update\_order.py deployed in Lambda}

\section{Part 8: Verification After Fix}
\begin{enumerate}
  \item Re-ran the same Python exploit script after deploying the fix.
  \item The API now returns \texttt{\{"status":"err","msg":"order already paid or in progress"\}}.
  \item Normal order placement and updates before payment still work as expected.
\end{enumerate}
\evidencebox{../evidence/lesson-08/image copy 4.png}{Post-fix: update attempt on paid order returns error}

\section[Part 9: Structured Operation and Security Analysis]{Part 9: Structured Operation and\\Security Analysis}
\subsection*{Table A}
\begin{table}[H]
  \small
  \centering
  \begin{tabularx}{\textwidth}{|Y|Y|Y|Y|}
    \hline
    \textbf{Intended rule(s)} & \textbf{Artifacts used to infer the rule} & \textbf{Normal evidence} & \textbf{Exploit evidence} \\
    \hline
    An order's items must be immutable once payment processing begins (status $>$ 100). &
    Order state machine in code comments and normal UI workflow. &
    Users can modify cart only before clicking ``Pay''. &
    Attacker modifies items after payment with a direct API call. \\
    \hline
  \end{tabularx}
\end{table}

\subsection*{Table B}
\begin{table}[H]
  \footnotesize
  \centering
  \begin{tabularx}{\textwidth}{|Y|Y|Y|Y|Y|}
    \hline
    \textbf{Why the exploit is a deviation} & \textbf{Deviation class} & \textbf{Fix applied} & \textbf{Post-fix verification} & \textbf{Optional latency / timing note} \\
    \hline
    Order contents were changed after billing started, violating order immutability. &
    Logic flaw / Race condition. &
    Changed status threshold from $>110$ to $>100$ in update\_order.py. &
    Update requests on non-open orders now return an error. &
    Exploit must be sent during the billing window (status 100$\rightarrow$120). \\
    \hline
  \end{tabularx}
\end{table}

\section{Part 10: Takeaway / Lessons Learned}
Business logic flaws are among the hardest vulnerabilities to detect with automated tools because they require understanding the intended workflow. Status-machine transitions must be strictly enforced server-side, and boundary conditions (off-by-one errors in status checks) must be carefully reviewed.

\clearpage

%%%%%%%%%%%%%%%%%%%%%%%%%%%%%%%%%%%%%%%%%%%%%%%%%%%%%%%%%%
\chapter{Lesson 09: Vulnerable Dependencies}
%%%%%%%%%%%%%%%%%%%%%%%%%%%%%%%%%%%%%%%%%%%%%%%%%%%%%%%%%%

\section{Part 1: Goal and Vulnerability Summary}
The objective is to identify and remediate a known-vulnerable third-party dependency (\texttt{node-serialize} v0.0.4) used by the \texttt{DVSA-ORDER-MANAGER} Lambda. This package has a publicly disclosed critical RCE vulnerability (CVE-2017-5941, CVSS 9.8) and should never be used in production.

\section{Part 2: Why This Works / Root Cause}
The application's \texttt{package.json} includes \texttt{node-serialize} version 0.0.4 as a dependency. This package uses JavaScript's \texttt{eval()} internally when deserializing objects, enabling arbitrary code execution. The developers chose this library without auditing it for known CVEs.

\section{Part 3: Environment and Setup}
\begin{itemize}
  \item Lambda function: \texttt{DVSA-ORDER-MANAGER}.
  \item AWS services involved: Lambda.
  \item Relevant file: \texttt{package.json} in the Lambda deployment package.
  \item User or test context: Security auditor reviewing dependencies.
  \item Tools used: AWS Lambda Console (code editor), NVD (NIST National Vulnerability Database).
\end{itemize}

\section{Part 4: Reproduction Steps}
\begin{enumerate}
  \item Opened the \texttt{DVSA-ORDER-MANAGER} Lambda function in the AWS Console.
  \item Inspected \texttt{package.json}---found \texttt{"node-serialize": "0.0.4"} listed.
  \item Searched NVD for ``node-serialize''---found CVE-2017-5941 with CVSS 9.8 (Critical).
  \item The CVE confirms that \texttt{unserialize()} enables arbitrary code execution via crafted payloads.
\end{enumerate}

\section{Part 5: Evidence and Proof}
\begin{itemize}
  \item Lambda Console showing \texttt{package.json} with the vulnerable dependency.
  \item NVD page for CVE-2017-5941 showing the critical severity rating.
\end{itemize}
\evidencebox{../evidence/lesson-09/package_problem.png}{package.json showing node-serialize 0.0.4 dependency}
\evidencebox{../evidence/lesson-09/public_proof.png}{CVE-2017-5941: CVSS 9.8 Critical---node-serialize RCE}

\section{Part 6: Fix Strategy / Probable Mitigation}
\begin{itemize}
  \item Fix location: \texttt{package.json} in the \texttt{DVSA-ORDER-MANAGER} Lambda deployment package.
  \item Security property restored: Elimination of known-vulnerable code paths from the runtime.
  \item Why this addresses root cause: Removing the dependency entirely eliminates the \texttt{eval()}-based deserialization vector. The code was already updated (Lesson 01) to use \texttt{JSON.parse()}, so the library is no longer needed.
\end{itemize}

\section{Part 7: Code / Config Changes}
Removed \texttt{node-serialize} from \texttt{package.json}:
\begin{lstlisting}[style=reportcode]
// BEFORE:
"dependencies": {
  "node-jose": "2.2.0",
  "node-serialize": "0.0.4"
}

// AFTER:
"dependencies": {
  "node-jose": "2.2.0"
}
\end{lstlisting}
Redeployed the Lambda function.

\evidencebox{../evidence/lesson-09/post_fix.png}{Post-fix: package.json with node-serialize removed, Lambda deployed successfully}

\section{Part 8: Verification After Fix}
\begin{enumerate}
  \item Confirmed \texttt{package.json} no longer lists \texttt{node-serialize}.
  \item Lambda deployed successfully without the vulnerable package.
  \item Normal order operations still function correctly (the code uses \texttt{JSON.parse()} and \texttt{node-jose} only).
\end{enumerate}

\section[Part 9: Structured Operation and Security Analysis]{Part 9: Structured Operation and\\Security Analysis}
\subsection*{Table A}
\begin{table}[H]
  \small
  \centering
  \begin{tabularx}{\textwidth}{|Y|Y|Y|Y|}
    \hline
    \textbf{Intended rule(s)} & \textbf{Artifacts used to infer the rule} & \textbf{Normal evidence} & \textbf{Exploit evidence} \\
    \hline
    Application dependencies must not contain known critical CVEs. &
    NVD database and package.json audit. &
    A clean dependency tree with no known vulnerabilities. &
    node-serialize 0.0.4 is present with CVE-2017-5941 (CVSS 9.8). \\
    \hline
  \end{tabularx}
\end{table}

\subsection*{Table B}
\begin{table}[H]
  \footnotesize
  \centering
  \begin{tabularx}{\textwidth}{|Y|Y|Y|Y|Y|}
    \hline
    \textbf{Why the exploit is a deviation} & \textbf{Deviation class} & \textbf{Fix applied} & \textbf{Post-fix verification} & \textbf{Optional latency / timing note} \\
    \hline
    A dependency with a publicly known critical RCE is shipped in production. &
    Vulnerable dependency / Supply-chain risk. &
    Removed node-serialize from package.json and redeployed. &
    Lambda runs successfully without the vulnerable package. &
    N/A \\
    \hline
  \end{tabularx}
\end{table}

\section{Part 10: Takeaway / Lessons Learned}
Dependencies must be audited regularly using tools like \texttt{npm audit}, Snyk, or Dependabot. A single vulnerable package can undermine an otherwise secure application. The principle of minimal dependencies---only include what is actively used---reduces attack surface significantly.

\clearpage

%%%%%%%%%%%%%%%%%%%%%%%%%%%%%%%%%%%%%%%%%%%%%%%%%%%%%%%%%%
\chapter{Lesson 10: Unhandled Exceptions}
%%%%%%%%%%%%%%%%%%%%%%%%%%%%%%%%%%%%%%%%%%%%%%%%%%%%%%%%%%

\section{Part 1: Goal and Vulnerability Summary}
The objective is to exploit unhandled exceptions in the \texttt{DVSA-ORDER-GET} Lambda function that leak internal implementation details (file paths, stack traces, variable names) to the client. When the function receives a malformed request missing required fields, it throws an unhandled \texttt{KeyError} that is returned verbatim to the caller.

\section{Part 2: Why This Works / Root Cause}
The original \texttt{get\_order.py} code directly accessed \texttt{event["orderId"]} without any try/except block or input validation. When the key is missing, Python raises a \texttt{KeyError} with a full traceback including the file path (\texttt{/var/task/get\_order.py}), line number, and variable names. This traceback was passed back through the Lambda response to the API client.

\section{Part 3: Environment and Setup}
\begin{itemize}
  \item API endpoint: \url{https://shfz8h7a28.execute-api.us-east-1.amazonaws.com/Stage/order}
  \item AWS services involved: API Gateway, Lambda, CloudWatch.
  \item Relevant function: \texttt{DVSA-ORDER-GET} (\texttt{get\_order.py}).
  \item User or test context: Authenticated user sending a malformed request.
  \item Tools used: Browser DevTools, \texttt{curl}, CloudWatch Logs.
\end{itemize}

\section{Part 4: Reproduction Steps}
\begin{enumerate}
  \item Sent a request to the \texttt{/order} endpoint with \texttt{\{"action":"get"\}} (missing the \texttt{order-id} field).
  \item The API returned a response containing the full Python traceback: \texttt{KeyError: 'orderId'} with file path and line number.
  \item This information disclosure reveals internal implementation details to an attacker.
\end{enumerate}

\section{Part 5: Evidence and Proof}
\begin{itemize}
  \item Browser DevTools showing the error response with stack trace and internal file paths.
  \item CloudWatch logs showing the unhandled KeyError with full traceback.
\end{itemize}
\evidencebox{../evidence/lesson-10/unhandeld exception.png}{Browser DevTools: error response leaks file path and stack trace}
\evidencebox{../evidence/lesson-10/log proof.png}{CloudWatch: KeyError traceback with internal file path exposed}

\section{Part 6: Fix Strategy / Probable Mitigation}
\begin{itemize}
  \item Fix location: \texttt{DVSA-ORDER-GET} Lambda, file \texttt{get\_order.py}.
  \item Security property restored: Safe error handling---internal details never reach the client.
  \item Why this addresses root cause: Wrapping the handler in try/except with a generic client-facing error message ensures stack traces are only logged server-side, never returned to callers.
\end{itemize}

\section{Part 7: Code / Config Changes}
Rewrote \texttt{get\_order.py} with proper error handling:
\begin{lstlisting}[style=reportcode]
def safe_error(message="bad request"):
    """Client-safe error response."""
    return {"status": "err", "msg": message}

def lambda_handler(event, context):
    try:
        order_id = get_required_field(event,
            ["orderId", "order-id", "order_id"])
        if not order_id:
            return safe_error("bad request")
        # ... normal logic ...
    except Exception as e:
        print("ERROR:", repr(e))  # log internally
        return safe_error("internal error")
\end{lstlisting}

\section{Part 8: Verification After Fix}
\begin{enumerate}
  \item Sent the same malformed request: \texttt{\{"action":"get"\}} with no order-id.
  \item The API now returns a generic \texttt{\{"status":"err","msg":"bad request"\}} with no stack trace.
  \item CloudWatch logs show the error is logged internally with ``ERROR: Missing order ID'' but not exposed to the client.
\end{enumerate}
\evidencebox{../evidence/lesson-10/post fix.png}{Post-fix: API returns generic ``bad request'' error, no stack trace}
\evidencebox{../evidence/lesson-10/post fix log.png}{Post-fix: CloudWatch shows safe internal logging only}

\section[Part 9: Structured Operation and Security Analysis]{Part 9: Structured Operation and\\Security Analysis}
\subsection*{Table A}
\begin{table}[H]
  \small
  \centering
  \begin{tabularx}{\textwidth}{|Y|Y|Y|Y|}
    \hline
    \textbf{Intended rule(s)} & \textbf{Artifacts used to infer the rule} & \textbf{Normal evidence} & \textbf{Exploit evidence} \\
    \hline
    Error responses must never expose internal implementation details (file paths, stack traces, variable names). &
    Secure coding best practices and OWASP guidelines. &
    The API returns generic error messages for invalid input. &
    Malformed request returns full Python traceback with file path and line numbers. \\
    \hline
  \end{tabularx}
\end{table}

\subsection*{Table B}
\begin{table}[H]
  \footnotesize
  \centering
  \begin{tabularx}{\textwidth}{|Y|Y|Y|Y|Y|}
    \hline
    \textbf{Why the exploit is a deviation} & \textbf{Deviation class} & \textbf{Fix applied} & \textbf{Post-fix verification} & \textbf{Optional latency / timing note} \\
    \hline
    Internal file paths and code structure are leaked to external users. &
    Unsafe error handling / Information disclosure. &
    Added try/except with a generic safe\_error() response function. &
    Same malformed request now returns only ``bad request'' with no internals. &
    N/A \\
    \hline
  \end{tabularx}
\end{table}

\section{Part 10: Takeaway / Lessons Learned}
Unhandled exceptions are a common vector for information disclosure in serverless applications. Every Lambda handler must wrap its logic in a try/except block that returns sanitized error messages to clients while logging full details server-side. This prevents attackers from fingerprinting the technology stack and discovering internal code paths.

\clearpage

%%%%%%%%%%%%%%%%%%%%%%%%%%%%%%%%%%%%%%%%%%%%%%%%%%%%%%%%%%

\chapter*{Conclusion}
\addcontentsline{toc}{chapter}{Conclusion}

This project demonstrated that serverless architectures, while abstracting infrastructure management, still require rigorous security engineering across multiple layers. Through 10 lessons on the DVSA platform, we identified and remediated vulnerabilities spanning insecure deserialization (Lesson 1), broken authentication (Lesson 2), sensitive data exposure (Lesson 3), insecure cloud configuration (Lesson 4), broken access control (Lesson 5), denial of service (Lesson 6), over-privileged IAM roles (Lesson 7), business logic flaws (Lesson 8), vulnerable dependencies (Lesson 9), and unsafe error handling (Lesson 10).

Key remediation themes:
\begin{itemize}
  \item \textbf{Defense in depth:} Security controls must exist at multiple layers---API Gateway, Lambda code, IAM policies, and cloud resource configuration.
  \item \textbf{Least privilege:} IAM roles must be scoped to only the actions and resources each function actually uses.
  \item \textbf{Input validation:} All user input must be treated as untrusted; use safe parsers and validate before processing.
  \item \textbf{Dependency hygiene:} Regularly audit third-party packages for known CVEs and remove unused dependencies.
  \item \textbf{Safe error handling:} Never expose internal details to clients; log verbosely server-side but respond generically.
\end{itemize}

\end{document}
